For each integer $n\ge 1$, define $f(n)$ to be the minimal integer $m\ge 1$ such that there exists some $k\ge 1$ with
\[
n = d_1(m)+d_2(m)+\cdots + d_k(m),
\]
where $1=d_1(m)<d_2(m)<\cdots<d_{\tau(m)}(m)=m$ are the positive divisors of $m$ in increasing order.
(Thus the right-hand side is the sum of the $k$ smallest divisors of $m$.)
The problem text notes that $f(2)$ and $f(5)$ are undefined (no such $m$).
Questions:
\begin{itemize}
\item Is $f(n)=o(n)$ as $n\to\infty$?
\item If not, is $f(n)=o(n)$ for almost all $n$, and is $\limsup_{n\to\infty} f(n)/n=\infty$?
\end{itemize}
